\subsection{PROJECTIVE LINEAR MAPPINGS}

Let V, \(V'\) be two left vector spaces over a field K, f a linear mapping of V into \(V'\) and \(N = \bar{f}^{-1}(0)\) its kernel. It is immediate that the image under f of a line (passing through 0) in V not contained in N is a line (passing through 0) in \(V'\) ; hence, on passing to the quotients, f defines a mapping g of \(\mathbf{P}(\mathbf{V}) - \mathbf{P}(\mathbf{N})\) into \(\mathbf{P}(\mathbf{V}')\) . Such a mapping is called a projective linear mapping (or, simply, a projective mapping); although it is defined on \(\mathbf{P}(\mathbf{V}) - \mathbf{P}(\mathbf{N})\) and not on \(\mathbf{P}(\mathbf{V})\) (when \(N \neq \{0\}\) ), we shall say by an abuse of language that g is a projective mapping of \(\mathbf{P}(\mathbf{V})\) into \(\mathbf{P}(\mathbf{V}')\) . The projective linear variety \(\mathbf{P}(\mathbf{N})\) , where g is not defined, is called the centre of g.

Note that, when \(g\) is defined on the whole of \(\mathbf{P}(\mathbf{V})\) (that is when \(\mathbf{N} = \{0\}\) ), \(g\) is an injection of \(\mathbf{P}(\mathbf{V})\) into \(\mathbf{P}(\mathbf{V}')\) .

When bases \((a_{\lambda})_{\lambda\in\mathbf{L}}, (b_{\mu})_{\mu\in\mathbf{M}}\) are given in V and \(V'\) respectively, a projective mapping of \(\mathbf{P}(\mathbf{V})\) into \(\mathbf{P}(\mathbf{V}')\) maps a point of \(\mathbf{P}(\mathbf{V})\) with homogeneous coordinates \(\xi_{\lambda}\) ( \(\lambda \in \mathbf{L}\) ) to a point of \(\mathbf{P}(\mathbf{V}')\) with a system of homogeneous coordinates \(\eta_{\mu}\) ( \(\mu \in \mathbf{M}\) ) of the form

\[
\eta_ {\mu} = \sum_ {\lambda \in L} \xi_ {\lambda} \alpha_ {\lambda \mu} \quad (\alpha_ {\lambda \mu} \in K).\tag{6}
\]

The centre of \(g\) is the linear variety defined by the equations

\[
\sum_ {\lambda \in L} \xi_ {\lambda} \alpha_ {\lambda \mu} = 0 \quad (\mu \in M).
\]

If C is the centre of g and M is a linear variety of \(\mathbf{P}(\mathbf{V})\) , the image under g of \(\mathbf{M} - (\mathbf{M} \cap \mathbf{C})\) is a linear variety of \(\mathbf{P}(\mathbf{V}')\) denoted (by an abuse of language) by \(g(\mathbf{M})\) . Then

\[
\dim g (\mathbf {M}) + \dim (\mathbf {M} \cap \mathbf {C}) + 1 = \dim \mathbf {M}\tag{7}
\]

(§ 7, no. 4, formula (12)). If \(\mathbf{M}'\) is a linear variety of \(\mathbf{P}(\mathbf{V}')\) , \(\bar{g}^{-1}(\mathbf{M}') \cup \mathbf{C}\) is a linear variety of \(\mathbf{P}(\mathbf{V})\) and

\[
\dim (\bar {g} ^ {- 1} (\mathbf {M} ^ {\prime}) \cup \mathbf {C}) = \dim \mathbf {C} + \dim (\mathbf {M} ^ {\prime} \cap g (\mathbf {P} (\mathbf {V}))) + 1.\tag{8}
\]

It is said, by an abuse of language, that \(\bar{g}^{-1}(\mathbf{M}') \cup \mathbf{C}\) is the inverse image of \(\mathbf{M}'\) under \(g\) .

As the values taken by a linear mapping on a basis \((e_{t})\) of V can be chosen arbitrarily in \(V'\) , it is seen that there exists a projective linear mapping of \(\mathbf{P}(\mathbf{V})\) into \(\mathbf{P}(\mathbf{V}')\) taking arbitrary values at the points \(\pi(e_{t})\) . But (even when g is everywhere defined) giving \(g(\pi(e_{t}))\) does not determine g uniquely (Exercise 10).

The composition of two projective mappings which are bijections is a projective mapping; so is the inverse mapping of such a bijection. The bijective projective mappings of a projective space \(\mathbf{P}(\mathbf{V})\) onto itself thus form a group, called the projective group of \(\mathbf{P}(\mathbf{V})\) and denoted by \(\mathbf{PGL}(\mathbf{V})\) ; we write \(\mathbf{PGL}_{n}(\mathbf{K})\) or \(\mathbf{PGL}(n,\mathbf{K})\) instead of \(\mathbf{PGL}(\mathbf{K}_{s}^{n})\) .

Remark. In a projective space \(\mathbf{P}(\mathbf{V})\) over a field \(\mathbf{K}\) , let \(\mathbf{H} = \mathbf{P}(\mathbf{W})\) be a hyperplane. There exists a bijective linear mapping \(f\) of \(\mathbf{V}\) onto \(\mathbf{K}_s \times \mathbf{W}\) such that \(f(\mathbf{W}) = \mathbf{W}\) ; let \(g\) be the projective mapping obtained from \(f\) by passing to the quotients. It has been seen (no. 8) that the complement of \(\mathbf{P}(\mathbf{W})\) in \(\mathbf{P}(\mathbf{K}_s \times \mathbf{W})\) can be identified with an affine space whose translation space is \(\mathbf{W}\) . When \(\mathbf{P}(\mathbf{V})\) is identified with \(\mathbf{P}(\mathbf{K}_s \times \mathbf{W})\) by means of \(g\) , it is said that \(\mathbf{H}\) has been taken as hyperplane at infinity in \(\mathbf{P}(\mathbf{V})\) ; the complement of \(\mathbf{H}\) in \(\mathbf{P}(\mathbf{V})\) is then identified with an affine space whose translation space is \(\mathbf{W}\) .

